\section{Complete ten-integer definition and transfer proof}
\label{sec:ten-supplement}

This supplement specifies the second delivered polynomial without any
implicit auxiliary unknowns.  Its ten independent unknowns are
\[
 (f,g,h,k,l,w,x,y,z,m)\in\Z^{10}.
\]
Every other symbol in this section abbreviates an integer polynomial in
these ten variables.  Symbols used for the nine-variable construction
are reset locally.  The parent polynomials $p_j$, their fixed weighted
coefficients $c_{j,\alpha}$, and either one of the two weight vectors
are exactly those in Sections~\ref{sec:parent}--\ref{sec:sparse}.
In particular,
\[
 W=\max_iw_i,\quad h_j=j(2W+1)-1,\quad
 H_*=44W+21,\quad K_*=46W+21,
\]
and $\mathcal D$ is the fixed univariate polynomial in
\eqref{eq:D}.  These are fixed exponents and fixed coefficients,
not further unknowns.

\subsection{The two-input packed source}

Set the fixed integer $\beta_{10}=2^{2140030}$ and define
\begin{align*}
 b&=1+3f,& B_{10}&=\beta_{10}b^2,
 &\ell&=\sum_{i=1}^{31}B_{10}^{w_i},\\
 N_0&=B_{10}^{W+1},&N_1&=4B_{10}^{K_*+1},&N&=N_0N_1,\\
 M_c&=N_0-1-(b-1)\ell,
 &Q_c&=(B_{10}-2)\sum_{j=1}^{22}B_{10}^{h_j},
 &T_c&=M_c+Q_cN_0,\\
 J_c&=2(1+g)^2\mathcal D(B_{10})+
                B_{10}\sum_{q=0}^{K_*}B_{10}^{q},
 &S_c&=g+J_cN_0.
\end{align*}
Finally put
\[
 X=(N-1)\bigl((S_c+T_c+1)N+T_c+1\bigr),\qquad
 Y=N^2,\qquad C_b=bB_{10}^{W}.
\]
Thus $N=4B_{10}^{47W+23}$ and
$C_b=2^{2140030W}b^{2W+1}$.
All sums have fixed finite limits.  The released circuit evaluates them
by shared powers and geometric-sum identities.

\subsection{The retained Lucas and norm registers}

Define, in the displayed order,
\begin{align*}
 L&=lY,&U&=2LX,&V&=4gwY,&W_b&=bw,\\
 A&=U(V+1),&B_1&=2X+1,&C&=B_1+(A-2)h,\\
 P_n&=U^2V,&K&=X+1+k(P_n-2),\\
 S&=2A-5,&T&=3W_bC-2(W_b^2-1),\\
 R_2&=(P_n^2-4)K^2+4.
\end{align*}
The margin and Lucas registers are
\begin{align*}
 \Delta_b&=8C_b^3gK^2-
       g^2\bigl(32(C-KL)^2C_b^3+g^2K^2\bigr),\\
 O&=f^2l^2\Delta_b,&\xi&=zx,\\
 D&=(A^2-4)C^2+4,&E&=C^2D\xi,\\
 F&=4(A^2-4)E^2+1,
 &G^*&=2(A+2)(A-2)^2E^2-1,\\
 \Omega_{10}&=16E,
 &H^*&=C+(B_1-C)F+\Omega_{10}Fy,\\
 I^*&=((G^*)^2-1)(H^*)^2+1.
\end{align*}
The homogeneous quintic merger and the three radicands are
\begin{align*}
 \Sigma&=\xi S^5,&r&=T+Sy,\\
 \Theta_{10}&=(r^3-rS^2-S^3)(r^2+23S^2),\\
 r_0&=b\bigl(1+(3O-4)b^2z^2\bigr),&
 r_1&=R_2,&r_2&=DFI^*,\\
 v_1&=r_0,&v_2&=v_1(P_n^2+K^2+3),\\
 Z_n&=\Sigma m-\Theta_{10},\\
 a_0&=\Sigma^2r_0,&a_1&=\Sigma^2v_1^2r_1,
 &a_2&=\Sigma^2v_2^2r_2.
\end{align*}
Set $V_n=Z_n^2+a_2-a_0-a_1$.  The complete equation is
\begin{equation}
 \boxed{\quad
 F_{10}:=(V_n^2+4Z_n^2a_2-4a_0a_1)^2
              -16V_n^2Z_n^2a_2=0.\quad}
 \label{eq:supplement-F10}
\end{equation}
After the definitions are substituted, this is a single polynomial over
$\Z$ in precisely the ten displayed integer unknowns.  The files
\texttt{sparse220\_ten\_integer.json} and
\texttt{sparse239\_ten\_integer.json} specify the corresponding expanded
binary arithmetic schedules and bind every nonliteral leaf to its exact
materialized coefficient.

\subsection{The precise inherited source interface}

Here is the interface of package 40's source theorem used in the
composition~\cite{source40}.  Let $H_0$ bound the ordinary coefficients
of the shifted parent rows.  One may take
\[
 H_0=8(1+\beta+K_0+M_C+M_F)^2,
 \qquad \beta_{10}>4\cdot33^2H_0.
\]
The displayed inequality follows from the fixed port bit lengths; the
exact coefficient payloads also verify it directly.  Then:
\begin{enumerate}[label=\textup{(\roman*)}]
\item If the $p_j$ have a common zero in $\N^{31}$, then for every
sufficiently large positive integer $f$ for which $b=1+3f$ is both a
square and a power of two, there is a $g$ with
\[
 b\leq g<C_b,\qquad Y\mid\binom{2X}{X}.
\]
\item If $f\ne0$, $b$ is a square and a power of two,
$0\leq g<2C_b$, and $Y\mid\binom{2X}{X}$, then the $p_j$ have a
common zero in $\N^{31}$.
\item Already under $f\ne0$, $b$ a square, and $0\leq g<2C_b$,
without the power-of-two assumption, one has
\[
 0\leq S_c,T_c<N,\quad3\mid X,\quad X>3b,\quad
 Y>\max\{b,256\}.
\]
\end{enumerate}

The new support-isolation property proves the exact same coefficient
identities and coefficient bound used in that theorem.  Distinct
encoded positions, positive leading coefficient 2, target spacing,
and nonnegative exponents have all been retained.  The proof of
\textup{(iii)} bounds $J_c$ for every integer in the broad permitted
interval before any digit recovery.  For \textup{(ii)}, the lower mask
recovers $g=\sum_i n_iB_{10}^{w_i}$ with $0\leq n_i<b$; the upper mask
tests all 22 centered coefficients, and the binary carry identity gives
the binomial divisibility.  For \textup{(i)}, choose arbitrarily large
$b=4^q>\max_i n_i$ and $f=(b-1)/3$, then use the same word.  The parent
zero is nonzero because its last row has constant coefficient 1; thus
$g\geq B_{10}\geq b$.  These observations prove that the source
interface transfers to the two supplied vectors.

Sun's retained argument uses this interface, $b=1+3f$, and
$C_b=2^c b^d$ with $c$ even and $d$ odd.  Here
$c=2140030W$ is even and $d=2W+1$ is odd.  Thus the new source has all
the required congruence and growth hypotheses.  The use is of the
parametric proofs of Propositions 6.1--6.2 and Theorem 4.1 of
Sun's 2026 paper, as composed and documented in packages 31 and
40~\cite{sun2026,source31,source40}.

\subsection{Why the unrestricted norm has exactly the required zeros}

The retained bridge has two relevant consequences.  A common parent
zero supplies positive witnesses with $A>2$ even, $B_1>0$ odd,
$b$ square, $b\leq g<C_b$, $C=u_{B_1}(A)$, $R_2$ square,
$S\mid T$, and $O>0$.  Conversely, for the formulas above,
\[
 b\text{ square},\quad O>0,\quad\xi\ne0,\quad
 DFI^*\text{ square},\quad R_2\text{ square},\quad S\mid T
 \tag{S1}
\]
imply a parent zero.  The reverse proof first obtains the source's
broad ranges from $O>0$; it then recovers the Lucas value and the
power-of-two base.  The full interpolation assertion says that, when
$C=u_{B_1}(A)$ under the positive hypotheses, for any positive $d$
one may choose $\xi>0$ divisible by $d$ such that \emph{every}
residue class modulo $\xi$ contains a $y$ making $DFI^*$ a square.

For completeness, the finalizer enforces exactly (S1) as follows.
The polynomial $\pi(t)=(t^3-t-1)(t^2+23)$ has no rational zero and has
a root modulo every positive integer.  Its cubic discriminant is
$-23$: at an odd prime other than 23, either the cubic has a simple
root or its irreducible Frobenius 3-cycle makes $-23$ a square; at 23
the cubic has simple root 3; at 2 one uses $-23\equiv1\pmod8$.
Lifting and CRT finish the local assertion.  As $S=2A-5$ is always
odd and nonzero, $\Theta_{10}\ne0$ always.  Moreover,
$\Sigma\mid\Theta_{10}$ implies $S\mid T$: if
$\nu_p(S)>\nu_p(T)$, the two homogeneous factors have valuations
$3\nu_p(T)$ and $2\nu_p(T)$, contradicting $S^5\mid\Theta_{10}$.
Conversely, from $S\mid T$ a root of $\pi$ modulo $|\xi|$ prescribes
a residue class for $y$ giving $\Sigma\mid\Theta_{10}$.

For $z\ne0$, the factors $b$ and $1+(3O-4)b^2z^2$ are coprime, and
the second cannot vanish.  Therefore $r_0$ square makes $b$ a signed
square; $b\equiv1\pmod3$ excludes the negative sign.  Its second
factor is then a square, which forces $O>0$.  Conversely positive
bridge witnesses have $O\geq g\geq b\geq4$.  Pell's equation for the
positive nonsquare $3O-4\equiv2\pmod3$, after taking a unit power
whose second coordinate is divisible by $b$, supplies $z>0$ making
$r_0$ square.

Finally, for $z\ne0$ one has $|r_0|>1$ and
\[
 \sqrt{|r_1|}\leq P_n^2+K^2+2,\qquad
 |v_1|>\sqrt{|r_0|},\qquad
 |v_2|>\sqrt{|r_0|}+|v_1|\sqrt{|r_1|}.
\]
These inequalities ensure that a rational signed sum
$\pm\sqrt{r_0}\pm v_1\sqrt{r_1}\pm v_2\sqrt{r_2}$ requires all
three radicands to be integer squares.  Indeed, a multiquadratic-field
automorphism changing a root changes the highest-index changed term
by more than the combined absolute values of all earlier changes.
The argument also covers negative radicands and dependent square
classes.

Equation~\eqref{eq:supplement-F10} is identically the product of its
eight signed-root factors.  If $\Sigma=0$, it equals
$\Theta_{10}^8\ne0$; hence any zero has $\xi=zx\ne0$.  Dividing a
vanishing factor by $\Sigma$ and using separation forces the three
squares and then $\Sigma\mid\Theta_{10}$.  The preceding mergers
give (S1), so the reverse bridge supplies a parent zero.  In the
other direction, choose the positive bridge witnesses, then the Pell
value $z>0$, then an interpolation modulus $\xi$ divisible by $z$,
then a $y$ in the quintic-prescribed residue class.  Put $x=\xi/z$.
All radicands are squares and $\Sigma\mid\Theta_{10}$, so an integer
$m$ makes one norm factor vanish.  This proves
\[
 (\exists\mathbf t\in\Z^{10})\ F_{10}(\mathbf t)=0
 \quad\Longleftrightarrow\quad
 (\exists\boldsymbol n\in\N^{31})\ \bigwedge_j p_j(\boldsymbol n)=0.
\]
Together with the inherited fixed-compiler theorem, classical negation
gives the asserted equivalence with $\Con$ in ZFC.  The argument
retains the cited ordinary proof dependencies and makes no additional
claim of end-to-end proof-assistant certification.

\subsection{Exact degree and supplied evaluators}

Writing $p=47W+23$, the source degrees are
$d_X=6p$, $d_Y=4p$, and $d_C=2W+1$.  The retained tail gives
\begin{align*}
 \deg Z_n&=5d_X+10d_Y+30,\\
 \deg a_0&=14d_X+28d_Y+3d_C+57,\\
 \deg a_1&=26d_X+48d_Y+6d_C+98,\\
 \deg a_2&=114d_X+224d_Y+6d_C+440.
\end{align*}
The last strictly dominates the degrees of $a_0,a_1,Z_n^2$.
Thus the top homogeneous part of $F_{10}$ is the fourth power of the
top homogeneous part of $a_2$, and
\[
 \deg F_{10}=297088W+147144.
\]
The $W=220$ evaluator uses 675 operations and has exact degree
65,506,504; the $W=239$ evaluator uses 662 operations and has exact
degree 71,151,176.  Independent direct leading-coefficient propagation
on these two exported DAGs gives nonzero residues 489390939 and
862326449 modulo 1,000,000,007, respectively, certifying attainment of
the propagated upper bounds.
